\section{Noncommutative and abstract Giry (co)monads}\label{sec:giry}
As a consequence of monadicity, we can infer crucial results regarding the coKleisli nature of $\SigmaCcpu$ and $\cSigmaCcpu$ by virtue of the adjoint lifting theorem~\cite[Theorem 4.5.6]{borceux94handbook2}.

We start with a preliminary result, which previously appeared in the context of pre-\cstar{}s~\cite[Proposition 4.18]{fritz2023involutive}.
For the sake of transparency, we offer a complete argument.

Recall that, given a \cstar{} $\A$, the \newterm{state space} $S(\A)$ is the set of states $\phi\colon \A \to \bbC$ equipped with the weakest topology making the evaluations
\[
	\ev_a \colon \xymatrix@R=0.1em{
	S(\A) \ar[r]
	& \bbC \\
	\phi  \ar@{|->}[r]
	& \phi(a)
	}
\]
continuous for all $a\in \A$.
It is well-known that $S(\A)$ is a compact Hausdorff space: for instance, this follows from the discussion of \cite[3.2.1]{pedersen2018automgroups}, which shows that $S(\A)$ is a compact subspace of $A^*$ equipped with the weak-$*$ topology.

\begin{proposition}\label{prop:state_space_leftadjoint}
	The forgetful functor $R_c\colon \cCalg \to \Calgcpu$ has a left adjoint given by\footnote{The following proof strategy was suggested to us by Tobias Fritz.}
	\[
		L_c\colon \xymatrix@R=.1em{
		\Calgcpu \ar[r]
		& \cCalg \\
		\A  \ar@{|->}[r]
		& C(S(\A))
		}
	\]
\end{proposition}
\begin{proof}
	Let $\CHaus$ be the category of compact Hausdorff spaces with continuous maps.
	For brevity, maps $\phi\colon \A \times X\to \bbC$, where $\A$ is a \cstar{} and $X$ is a compact Hausdorff space, are called \emph{state-continuous maps} if $\phi(-,x)\colon \A \to \bbC$ are states for all $x \in X$ and $\phi(a,-)\colon X\to \bbC$ are continuous for all $a\in \A$  (cf.\ the definition of Markov kernels, recalled on p.~\pageref{def:mkernel}); the set of these maps will be denoted by $\mathbf{SCM}(\A\times X,\bbC)$.

	The statement is obtained by the following chain of isomorphisms (whose naturality in $\A$ and $X$ is a direct check):
	\begin{align*}
		\cCalg (C(S(\A)),C(X)) & \cong \CHaus(X,S(\A))                                        \\
		                       & \cong \mathbf{SCM}(\A\times X,\bbC) \tag{$\spadesuit$} \\
		                       & \cong \Calgcpu(\A,C(X)) \tag{$\heartsuit$}
	\end{align*}
	While the first isomorphism is an immediate consequence of Gelfand duality, the remaining two require further explanation.
	We do so by exhibiting mutually inverse maps.
	\begin{itemize}
		\item For $(\spadesuit)$, we define
		      \[\begin{gathered}
			      \xymatrix@R=.1pt{
			      \CHaus(X,S(\A)) \ar[r]
			      & \mathbf{SCM}(\A\times X,\bbC) \\
			      \phi  \ar@{|->}[r]
			      & \bar{\phi}(a,x)\coloneqq \phi(x)(a)
			      } \\
			      \xymatrix@R=.1pt{
			      \mathbf{SCM}(\A\times X,\bbC)  \ar[r]
			      & \CHaus(X,S(\A)) \\
			      \psi  \ar@{|->}[r]
			      & \bar{\psi}(x)\coloneqq  \psi(-,x)
			      }
								\end{gathered}\]
		      These maps are well-defined: $\bar{\phi}$ is state-continuous because $\phi(x)$ is a state by definition of $S(\A)$ and $\phi(-)(a)=\ev_a\circ\, \phi$ is continuous because it is a composition of continuous maps. For the second one, by definition of $S(\A)$, continuity of $\bar{\psi}$ can be checked by postcomposition with $\ev_a$, and $\ev_a \circ \bar{\psi} = \psi(a,-)$ is continuous by assumption.
		      That these are inverses of each other is a direct check.

		\item For $(\heartsuit)$, we consider
		      \[\begin{gathered}
			      \xymatrix@R=.1pt{
										\mathbf{SCM}(\A\times X,\bbC)  \ar[r]
										& \Calgcpu(\A,C(X)) \\
										\psi  \ar@{|->}[r]
										& \widetilde{\psi}(a)\coloneqq  \psi(a,-)
										} \\
										\xymatrix@R=.1pt{
											\Calgcpu(\A,C(X)) \ar[r]
											& \mathbf{SCM}(\A\times X,\bbC) \\
											\phi  \ar@{|->}[r]
											& \widetilde{\phi}(a,x)\coloneqq \phi(a)(x)
											}
										\end{gathered}
		      \]
		      Concerning the well-definedness of $\widetilde{\psi}$, we have that $\psi(a,-)$ is continuous by definition, so $\widetilde{\psi}(a)$ indeed belongs to $C(X)$.
		      To show that $\widetilde{\psi}$ is a \cpu{} map, take any positive element $a\in \A$: since $\psi(-,x)$ is a state, $\psi(a,x)$ is nonnegative.
		      As this holds for all $x$, $\psi(a,-)$ is a positive element of $C(X)$.
		      A similar reasoning ensures that $\widetilde{\psi}(1)=1$, and therefore $\widetilde{\psi}$ is a \cpu{} map because its codomain is commutative (\cite[Theorem 3.9]{paulsen2002}).
		      Finally, $\widetilde{\phi}$ is a state-continuous map because $\widetilde{\phi}(a,-)$ is continuous by definition, and $\widetilde{\phi}(-,x)$ is a state for $\A$ because it is a composition of $\phi$ with the state $C(X)\owns f\mapsto f(x)\in \bbC$.
		      As above, these maps are inverses of one another by direct check.
		      \qedhere
	\end{itemize}
	%%%%OLD PROOF
	% We use the description given by universal morphisms.
	% Explicitly, we show that $R_c$ is a right adjoint by providing a map $\ev_{\A}\colon \A \to R_c(L_c(\A))$ such that for every $\phi\colon \A \to R_c(\C)$, there exists a unique $\widetilde{\phi}\colon L_c(\A) \to \C$ satisfying 
	% \[
	% \xymatrix{\A\ar[rr]^{\ev_{\A}}\ar[rd]_{\phi}&&R_c(L_c(\A))\ar[ld]^{R_c(\widetilde{\phi})}\\
	% &R_c(\C)&
	% }
	% \]
	% The map $\ev_{\A}$ is defined by $a\mapsto \ev_{a}$: it is positive unital by direct computation, and so it is completely positive because $L_c(A)$ is commutative (\cite[Theorem 3.9]{paulsen2002}).
	%
	% Let us now construct $\widetilde{\phi}$. 
	% Since $\C$ is commutative, by Gelfand duality there is an isomorphism $\C\cong C(X)$ for some compact Hausdorff space $X$, so we can assume $\C=C(X)$ without loss of generality. 
	% Given any \cpu{} map $\phi \colon \A \to R_c(C(X))$, we consider the map 
	% \[ 
	% \bar{\phi}\colon
	% \xymatrix@R=.1em{
	%  X \ar[r]
	%     & S(\A) \\
	%   p  \ar@{|->}[r]
	%     & \left(\bar{\phi}_p\colon a\mapsto \phi(a)(p)\right)
	%   }
	% \]
	% The fact that $\bar{\phi}_p$ is a state is immediately verified: whenever $a\ge 0$, then $\phi(a)$ is a nonnegative valued function, i.e. $\phi(a) \colon X\to [0,\infty)$, and therefore $\phi(a)(p)\ge 0$.
	% Continuity of $\bar{\phi}$ is also a direct check, since postcomposing $\bar{\phi}$ with $\ev_a$ yields a continuous map by definition: $\ev_a \bar{\phi}=\phi(a)$.
	% By Gelfand duality, $\bar{\phi}$ uniquely corresponds to a $*$-homomorphism $\widetilde{\phi}\colon L_c(\A)= C(S(\A))\to C(X)$. 
	% By construction, for every $a \in \A$,
	% \[
	% R_c(\widetilde{\phi})(\ev_{\A}(a)) = \widetilde{\phi}(\ev_{a}) = \ev_{a}\bar{\phi} = \phi(a).
	% \]
	%
	% We are left to show that $\widetilde{\phi}$ is unique. 
	% This follows by noting that any $*$-homomorphism $L_c(\A)\to \C$ is uniquely determined by its values on the family $\lbrace \ev_a: a\in \A\rbrace$ because the $*$-subalgebra generated by this family is dense in $L_c(\A)$ by the Stone--Weierstrass theorem (and $*$-homomorphisms are bounded). 
\end{proof}

\begin{theorem}[Noncommutative and abstract Giry (co)monads]\label{thm:comonad}
The following assertions hold.
	\begin{enumerate}
		\item\label{it:SigmaCcpu_coKleisli} The forgetful functor $R'\colon \SigmaC\to \SigmaCcpu$ has a left adjoint $L'$, and the induced comonad $L'R'$ on $\SigmaC$ has $\SigmaCcpu$ as its coKleisli category. This also holds when restricting to Pedersen--Baire envelopes.
		\item\label{it:cSigmaCcpu_coKleisli} Analogously, the forgetful $R'_{cc}\colon \cSigmaC\to \cSigmaCcpu$ has a left adjoint $L'_{cc}$, and the induced comonad $L'_{cc}R'_{cc}$ on $\cSigmaC$ has $\cSigmaCcpu$ as its coKleisli. This also holds when restricting to commutative Pedersen--Baire envelopes.
		\item\label{it:state_space} The forgetful functor $R'_c\colon \cSigmaC \to \SigmaCcpu$ has a left adjoint. This also restricts to Pedersen--Baire envelopes.
	\end{enumerate}
\end{theorem}
\begin{proof}
	We focus on Item~\ref{it:SigmaCcpu_coKleisli}.
	By \cref{thm:U_monadic,rem:monadic_cpu}, the forgetful functors $U\colon \SigmaC \to \Calg$ and $\bar{U}\colon \SigmaCcpu \to \Calgcpu$ are monadic.
	Moreover,
	\begin{itemize}
		\item The forgetful functor $R\colon \Calg \to \Calgcpu$ has a left adjoint $L\colon \Calgcpu \to \Calg$ by \cite[Proposition 8]{kirchberg1998Cstarfromoperator}.
		\item The forgetful functor $R'\colon \SigmaC \to \SigmaCcpu$ satisfies $\bar{U}\circ R' = R\circ U$ by immediate check (this composition is the obvious forgetful functor).
	\end{itemize}
	Putting it all together, we obtain the following diagram
	\[
		\vxy{
		\SigmaC\ar[r]^-{R'} \ar[d]^U & \SigmaCcpu \ar[d]^{\bar{U}}\\
		\Calg \ar@<.4em>[r]^-R\ar@{<-}@<-.4em>[r]_-L\ar@{}[r]|-\top & \Calgcpu}
	\]
	whose vertical arrows are monadic.
	By the adjoint lifting theorem~\cite[Theorem 4.5.6]{borceux94handbook2}, since $\SigmaC$ is cocomplete by \cref{prop:cocomplete}, we conclude that $R'$ has a left adjoint $L'$.
	The fact that the coKleisli category of the comonad $L' \circ R'$ is isomorphic to $\SigmaCcpu$ follows from the fact that $R'$ is bijective on objects: for all \cstars{}s $\A$ and $\B$,
	\begin{align*}
		\coKl(L'R') (\A,\B) & \cong \SigmaC(L'R'(\A) , \B)\cong \SigmaCcpu(R'(\A),R'(\B)) \\
		                    & \cong \SigmaCcpu(\A,\B).
	\end{align*}
	The restriction to Pedersen--Baire envelopes is straightforward, as the diagram of left adjoint functors commutes because the diagram of right adjoints does: for all $\A$ and $\B$ in $\SigmaC$,
	\begin{align*}
		\SigmaC(L'\bar{F} \A,\B) & \cong  \SigmaCcpu (\bar{F}\A, R'\B) \cong \Calgcpu (\A,\bar{U}R' \B)      \\
		                         & \cong \Calgcpu (\A, RU\B) \cong \Calg(L\A,U\B) \cong \SigmaC(FL \A , \B),
	\end{align*}
	and so by the Yoneda lemma $L'\bar{F} \A\cong FL \A$.

	The proof of Items~\ref{it:cSigmaCcpu_coKleisli} and \ref{it:state_space} are analogous.
	They are obtained by the commutative diagrams
	\[
		\vxy{
		\cSigmaC\ar@<.4em>[r]^-{R'_{cc}}\ar@{<-}@<-.4em>[r]_-{L'_{cc}}\ar@{}[r]|-\top \ar[d]^U & \cSigmaCcpu \ar[d]^{\bar{U}}\\
		\cCalg \ar@<.4em>[r]^-{R_{cc}}\ar@{<-}@<-.4em>[r]_-{L_{cc}}\ar@{}[r]|-\top & \cCalgcpu} \qquad \vxy{
		\cSigmaC\ar@<.4em>[r]^-{R'_c}\ar@{<-}@<-.4em>[r]_-{L'_c}\ar@{}[r]|-\top \ar[d]^U & \SigmaCcpu \ar[d]^{\bar{U}}\\
		\cCalg \ar@<.4em>[r]^-{R_c}\ar@{<-}@<-.4em>[r]_-{L_c}\ar@{}[r]|-\top & \Calgcpu}
	\]
	where $U$ and $\bar{U}$ in both diagrams are monadic as observed in \cref{rem:monadic_cpu}, the left adjoints of $R_{cc}$ and $R_c$ are given by the state space (see \cref{prop:state_space_leftadjoint}), and therefore $L'_{cc}$ and $L'_c$ are again obtained by the adjoint lifting theorem.\footnote{Cocompleteness of $\cSigmaC$ can be obtained by adapting our \cref{prop:cocomplete} above. Alternatively, this holds because $\cSigmaC$ is equivalent to the category of \bsigma{}s and $\sigma$-homomorphisms, which is known to be complete and cocomplete (\cite[Remark 2.1.12]{fritz2025categories}). }
	It is worth stressing that while $R'_{cc}$ is bijective on objects, $R'_c$ is not: for this reason, $\cSigmaCcpu$ is a coKleisli category over $\cSigmaC$, whereas the same does not hold for $\SigmaCcpu$.
	The restriction to (commutative) Pedersen--Baire envelopes holds as before.
\end{proof}
	The name given to the previous result is not merely suggestive.
	We show in \cref{thm:giry_restriction} that, under the equivalence between Baire measurable spaces and commutative Pedersen--Baire envelopes, the comonad of \cref{thm:comonad} \emph{is} the Giry monad.
	The remainder of this section makes this precise, and may be skipped by the reader willing to take the name on trust.

The Giry monad, first appeared in \cite{giry82probability}, is a monad $\mathcal{G}$ acting on $\Meas$, the category of measurable spaces with measurable functions.
Given a measurable space $X$, $\mathcal{G}(X)$ is given by the set of probability measures on $X$, equipped with the smallest $\sigma$-algebra making
\[
	\ev_U \colon \xymatrix@R=0pt{
	\mathcal{G}(X) \ar[r]
	& \bbC \\
	p  \ar@{|->}[r]
	& p(U)
	}
\]
measurable for all measurable subsets $U\subseteq X$ (cf.\ the state space of a \cstar{}).

The Kleisli category associated to $\mathcal{G}$ is $\Stoch$, the category of measurable spaces with Markov kernels (\cite[Section I.4]{giry82probability}).
As discussed in \cref{sec:preliminaries}, \scpu{} maps between commutative \cstars{}s may be viewed as ``abstract Markov kernels'', and similarly, one may view \scpu{} maps between \cstars{}s as ``noncommutative Markov kernels'', thus establishing a connection with the previous statement.
For the sake of convenience, we recall that a \newterm{Markov kernel} from $(X,\Sigma_X)$ to $(Y,\Sigma_Y)$ is a map $\phi\colon \Sigma_Y \times X \to [0,1]$ such that $\phi(U,-)\colon X\to [0,1]$ is measurable for all $U\in \Sigma_Y$ and $\phi(-,x)\colon \Sigma_Y\to [0,1]$ is a probability measure for all $x\in X$.\label{def:mkernel}

As argued in \cite{jamneshan23foundational}, a natural framework for measure theory that avoids second-countability constraints, unlike standard Borel spaces, is provided by \newterm{Baire measurable spaces}. These are measurable spaces $(X,\Sigma_X)$ where $\Sigma_X$ is the Baire $\sigma$-algebra of a compact Hausdorff topology on $X$.
We now seek to show that the Giry monad, when restricted to Baire measurable spaces, in fact coincides with the comonad structure obtained in \cref{thm:comonad}. 

\begin{theorem}\label{thm:giry_restriction}
	Let us consider:
	\begin{itemize}
		\item The category $\BaireMeas$ of Baire measurable spaces with measurable functions.
		\item The category $\cPBenv$ of commutative Pedersen--Baire envelopes with \starshom{}s.
	\end{itemize}
	Then the Giry monad restricts to $\BaireMeas$, i.e.\ $\mathcal{G}(X)$ is Baire measurable whenever $X$ is, and the following diagram commutes up to a natural isomorphism
	\[
		\vxy{
		\BaireMeas^{\op}\ar[r]^-{\mathcal{G}}\ar[d]^{\Linf} & \BaireMeas^{\op} \ar[d]^{\Linf}\\
		\cPBenv \ar[r]^-{\mathcal{S}'} & \cPBenv}
	\]
	where $\mathcal{S}'$ denotes the (restricted) comonad obtained by the forgetful functor $R'_{cc}\colon \cSigmaC \to \cSigmaCcpu$ and its adjoint.

\end{theorem}
The vertical arrows in the statement are equivalences \cite[Corollary 4.2.9]{fritz2025categories}.
\begin{proof}
	The fact that $\mathcal{G}(X)$ is Baire measurable whenever $X$ is follows from \cite[Corollary A.4]{shirazi2026commutativity}.
	To show the desired commutativity, we argue with the respective characterizations of the Kleisli and coKleisli categories.
	\begin{itemize}
		\item When restricted to $\BaireMeas$, the Kleisli category of the Giry monad is $\BaireStoch$, the category of Baire measurable spaces with Markov kernels. This follows from the fact that the Kleisli category of the Giry monad on $\Meas$ is $\Stoch$ (\cite[Section I.4]{giry82probability}).
		\item By \cref{thm:comonad}, the coKleisli category of $\mathcal{S}'$ is $\cpbenvcpu$, the category of commutative Pedersen--Baire envelopes with \scpu{} maps.
	\end{itemize}
	We claim that the diagram
	\begin{equation}\label{eq:Giry_extend_forget}
		\vxy{
			\BaireMeas^{\op}\ar[r]\ar[d]^{\Linf} & \BaireStoch^{\op} \ar[d]^{\Linf}\\
			\cPBenv \ar[r] & \cpbenvcpu}
	\end{equation}
	where the horizontal arrows are forgetful functors and the vertical arrows are equivalences \cite[Corollaries 4.2.9 and 4.3.8]{fritz2025categories}, is commutative.
	This is easily seen once all details have been made explicit.

	Let $\phi\colon X\to Y$ be an arbitrary measurable function in $\BaireMeas$. Via the forgetful $\BaireMeas\to \BaireStoch$, $\phi$ is sent to
	\[
		\bar{\phi} \colon \xymatrix@R=.1pt{
		\Sigma_Y\times X \ar[r]
		& [0,1] \\
		(U,x)
		\ar@{|->}[r]
		& \delta_{\phi(x)}(U),}
	\]
	where $\Sigma_Y$ is the $\sigma$-algebra of $Y$ and $\delta_{p}$ is the delta measure concentrated in $p$ (the unit of the Giry monad $X\to \mathcal{G}(X)$ is given by sending each point to its delta measure).
	The commutativity of \eqref{eq:Giry_extend_forget} is then equivalent to showing that $\Linf(\phi)= \Linf(\bar{\phi})$,
	which means that for every $f\colon Y\to \bbC$ and $x \in X$, $f(\phi(x))=\int_Y f(y) \bar{\phi}(\mathrm{d}y,x)$.

	But this holds by definition of $\bar{\phi}$:
	\[
		\int_Y f(y) \bar{\phi}(\mathrm{d}y,x)=\int_Y f(y) \delta_{\phi(x)}(\mathrm{d}y) =f(\phi(x)),
	\]
	so commutativity of \eqref{eq:Giry_extend_forget} is finally ensured.

	Since a square of left adjoints commutes whenever the corresponding square of right adjoints does, (as mentioned in the proof of \cref{thm:comonad}), and since $\Linf$ is an equivalence, we obtain the commutative diagram
	\[
		\vxy{
			\BaireStoch^{\op} \ar[r]\ar[d]^{\Linf} & \BaireMeas^{\op} \ar[d]^{\Linf}\\
			\cpbenvcpu \ar[r] & \cPBenv}
	\]
	where the horizontal arrows send $X$ to $\mathcal{G}(X)$ and $\A$ to $\mathcal{S}'(\A)$, respectively.
	We can now horizontally compose the two commutative diagrams to conclude the proof of the statement.
\end{proof}

\begin{remark}
	The categorically inclined reader may note that the proof above, making use of the uniqueness of left adjoints, directly implies that $\zeta\colon \Linf \mathcal{G}\Rightarrow \mathcal{S}'\Linf$ is compatible with the counit and comultiplication structures. Equivalently, the following diagrams of natural transformations between composite functors commute:
	\[\small
		\xymatrix{
		\Linf \circ\mathcal{G} \ar@{=>}[rr]^{\zeta} \ar@{=>}[dr]_{\Linf * \epsilon^{\mathcal{G}}} &&\mathcal{S}'\circ\Linf \ar@{=>}[dl]^{\epsilon^{\mathcal{S}'}* \Linf} \\
		& \Linf
		}% \qquad
		\xymatrix{
		\Linf \circ\mathcal{G} \ar@{=>}[rr]^{\zeta} \ar@{=>}[d]_{\Linf * \mu^{\mathcal{G}}}
		&& \mathcal{S}' \circ\Linf \ar@{=>}[d]^-{\mu^{\mathcal{S}'} * \Linf}
		& \\
		\Linf \circ\mathcal{G}\circ\mathcal{G} \ar@{=>}[r]_{\zeta *\mathcal{G}}
		& \mathcal{S}'\circ\Linf \circ\mathcal{G}\ar@{=>}[r]_{\mathcal{S}' *\zeta} & \mathcal{S}'\circ\mathcal{S}'\circ \Linf
		}
	\]
	where $\epsilon^{-} $ and $\mu^{-}$ denote the counit and comultiplication of the comonads $\mathcal{G}$ (on $\BaireMeas^{\op}$) and ${\mathcal{S}'}$.
\end{remark}


